% =========================================================================
% PORTADA OFICIAL
% =========================================================================
\includepdf[pages=-]{PORTADA.pdf}

% =========================================================================
% DEDICATORIA
% =========================================================================
\chapter*{Dedicatoria}
\addcontentsline{toc}{chapter}{Dedicatoria}
A nuestras familias y mentores que creyeron en este proyecto. A los usuarios y ahorradores del sector financiero popular en México, cuya tenacidad impulsa la búsqueda de herramientas cuantitativas más seguras y transparentes.

% =========================================================================
% AGRADECIMIENTOS
% =========================================================================
\chapter*{Agradecimientos}
\addcontentsline{toc}{chapter}{Agradecimientos}
Expresamos nuestra más sincera gratitud a nuestros directores de tesis, cuya guía, orientación y retroalimentación fueron fundamentales para la realización de este trabajo. Agradecemos también a los profesores de la Escuela Superior de Física y Matemáticas (ESFM) por su formación académica rigurosa, y a nuestras familias por su apoyo incondicional durante este proceso de investigación. Su confianza y motivación constante fueron esenciales para llevar esta tesis a término.

% =========================================================================
% RESUMEN
% =========================================================================
\chapter*{Resumen}
\addcontentsline{toc}{chapter}{Resumen}
Las Sociedades Financieras Populares (SOFIPOs) en México constituyen intermediarios financieros críticos para la inclusión de segmentos históricamente marginados del sistema bancario tradicional. Esta tesis desarrolla un sistema cuantitativo híbrido de alerta temprana que unifica econometría de series temporales de panel con algoritmos modernos de aprendizaje automático supervisado. Utilizando datos regulatorios oficiales de la CNBV correspondientes a 126 instituciones en operación durante el período 2015--2025 (1690 observaciones institución-trimestre depuradas), se construyó una arquitectura de procesamiento que permite análisis eficiente de paneles heterogéneos. Las variables macroeconómicas, es decir la TIIE a 28 días, la inflación, el tipo de cambio FIX, los Cetes a 28 días y el IGAE, se descargaron de la API del Sistema de Información Económica de Banco de México.

Tras la ingeniería de características (rezagos, medias móviles, ratios financieros), el panel de modelado quedó constituido por 1189 observaciones sobre 46 instituciones con historia suficiente. El modelo XGBoost, seleccionado mediante búsqueda de hiperparámetros con validación cruzada temporal, alcanzó un coeficiente de determinación ($R^2$) de \textit{0.6272} sobre el conjunto de prueba, con un Error Absoluto Medio (MAE) de \textit{3.16\%}, superando a LightGBM ($R^2=0.5756$) y resultando comparable a Random Forest ($R^2=0.5804$).

Para la tarea de clasificación de insolvencia crítica (IMOR futuro $>10\%$), se entrenó un ensamble apilado de Random Forest, XGBoost y LightGBM asistido por remuestreo híbrido SMOTE-ENN, obteniendo una exactitud de \textit{88.24\%}, precisión de \textit{69.49\%}, sensibilidad de \textit{80.39\%}, $F_1$ de \textit{0.7455}, área bajo la curva ROC de \textit{0.9227} y coeficiente de correlación de Matthews de \textit{0.6725}. Mediante valores de Shapley (SHAP) se identificó que el rezago del propio IMOR, el rezago del Nivel de Capitalización y el rezago del Resultado Neto son los predictores individuales de mayor peso, seguidos por el IGAE y el tipo de cambio entre las variables macroeconómicas. Todos los resultados fueron reproducidos de manera independiente a partir del código y los datos disponibles para este trabajo.

El marco teórico y el marco metodológico se desarrollan desde los conceptos más elementales hasta los más avanzados. Se explican la probabilidad y la estadística utilizadas, la estructura del estado de resultados y del balance de una SOFIPO y las razones por las que cada partida se emplea como variable, la inflación, las tasas de interés, el tipo de cambio y la actividad económica, la optimización matemática y computacional del código y cada paso del procesamiento de datos. Las gráficas generadas con Python se producen a una resolución de 350 puntos por pulgada. El trabajo documenta además los límites de su diseño de validación, entre ellos que la partición de entrenamiento y prueba no es cronológica y que la imputación de valores faltantes utiliza información futura.

\textbf{Palabras clave:} SOFIPO, riesgo crediticio, alerta temprana, XGBoost, SMOTE-ENN, SHAP, DuckDB, Polars, econometría de panel, insolvencia, estados financieros, optimización.

% =========================================================================
% ABSTRACT
% =========================================================================
\chapter*{Abstract}
\addcontentsline{toc}{chapter}{Abstract}
Popular Financial Companies (SOFIPOs) in Mexico are pivotal grassroots financial intermediaries serving populations historically excluded from commercial banking. This thesis develops a hybrid quantitative Early Warning System (EWS) that unifies panel time-series econometrics with modern supervised machine learning. Using official CNBV regulatory data for 126 operating institutions over 2015--2025 (1690 cleaned institution-quarter observations), a columnar processing architecture enables efficient analysis of heterogeneous panels. Macroeconomic variables, namely the 28-day TIIE, inflation, the FIX exchange rate, 28-day Cetes and the IGAE, were downloaded from the Banco de México Economic Information System API.

After feature engineering (lags, moving averages, financial ratios), the modeling panel comprised 1189 observations across 46 institutions with sufficient history. The XGBoost regressor, selected via time-series cross-validated hyperparameter search, achieved an test-set $R^2$ of \textit{0.6272} with a Mean Absolute Error of \textit{3.16\%}, outperforming LightGBM ($R^2=0.5756$) and matching Random Forest ($R^2=0.5804$).

For binary critical-insolvency classification (future NPL ratio $>10\%$), a Stacking Ensemble of Random Forest, XGBoost and LightGBM with SMOTE-ENN resampling attained \textit{88.24\%} accuracy, \textit{69.49\%} precision, \textit{80.39\%} recall, an $F_1$ of \textit{0.7455}, ROC-AUC of \textit{0.9227}, and a Matthews Correlation Coefficient of \textit{0.6725}. Shapley (SHAP) value decomposition identified the institution's own lagged NPL ratio, lagged capitalization level, and lagged net income as the strongest individual predictors, followed by the economic activity index and the exchange rate among macroeconomic drivers. All findings were independently reproduced from the code and data available for this study.

The theoretical and methodological frameworks are developed from the most elementary concepts to the most advanced ones. They explain the probability and statistics used, the structure of a SOFIPO's income statement and balance sheet and the reasons why each item is used as a variable, inflation, interest rates, the exchange rate and economic activity, the mathematical and computational optimization behind the code, and every data-processing step. All Python figures are produced at 350 dots per inch. The study also documents the limits of its validation design, including that the train-test split is not chronological and that imputation uses future information.

\textbf{Keywords:} Microfinance, credit risk, early warning systems, XGBoost, Stacking Ensemble, SHAP, panel econometrics, insolvency, non-performing loans, financial statements, optimization.

% =========================================================================
% LISTA DE ABREVIATURAS
% =========================================================================
\chapter*{Lista de Abreviaturas}
\addcontentsline{toc}{chapter}{Lista de Abreviaturas}
\begin{description}[leftmargin=3cm,style=nextline]
\item[AI / IA] Artificial Intelligence / Inteligencia Artificial
\item[ACF] Autocorrelation Function
\item[API] Application Programming Interface
\item[ARIMA] Autoregressive Integrated Moving Average
\item[Banxico] Banco de México
\item[CART] Classification and Regression Trees
\item[Cetes] Certificados de la Tesorería de la Federación
\item[CINIF] Consejo Mexicano de Normas de Información Financiera
\item[CNBV] Comisión Nacional Bancaria y de Valores
\item[dpi] Dots per inch / Puntos por pulgada
\item[EAD] Exposure at Default
\item[ENN] Edited Nearest Neighbors
\item[ETL] Extract, Transform, Load / Extracción, transformación y carga
\item[FIX] Tipo de cambio determinado por Banco de México (peso por dólar)
\item[GOSS] Gradient-based One-Side Sampling
\item[IGAE] Indicador Global de la Actividad Económica
\item[IMOR] Índice de Morosidad
\item[INEGI] Instituto Nacional de Estadística y Geografía
\item[INPC] Índice Nacional de Precios al Consumidor
\item[LGD] Loss Given Default
\item[LSTM] Long Short-Term Memory
\item[MAE] Mean Absolute Error
\item[MCC] Matthews Correlation Coefficient
\item[ML] Machine Learning
\item[NICAP] Nivel de Capitalización
\item[NIF] Normas de Información Financiera
\item[PACF] Partial Autocorrelation Function
\item[PD] Probability of Default
\item[PR-AUC] Area Under the Precision-Recall Curve
\item[RMSE] Root Mean Square Error
\item[RNN] Recurrent Neural Network
\item[ROA] Return on Assets
\item[ROC-AUC] Area Under the Receiver Operating Characteristic Curve
\item[ROE] Return on Equity
\item[R\textsuperscript{2}] Coeficiente de Determinación
\item[SARIMA] Seasonal ARIMA
\item[SHAP] SHapley Additive exPlanations
\item[SIE] Sistema de Información Económica
\item[SMOTE] Synthetic Minority Over-sampling Technique
\item[SOFIPO] Sociedad Financiera Popular
\item[SQL] Structured Query Language
\item[TIIE] Tasa de Interés Interbancaria de Equilibrio
\item[VaR] Value at Risk
\item[XGBoost] Extreme Gradient Boosting
\end{description}

\newpage
\tableofcontents
\newpage
\listoffigures
\newpage
\listoftables

